\documentclass[12pt]{article}
\usepackage[T1]{fontenc}
\usepackage[utf8]{inputenc}
\usepackage{lmodern}
\usepackage{textcomp}
\usepackage{enumitem}
\usepackage{geometry}
\geometry{margin=1in}
\begin{document}
\begin{center}
Answer Key - Economics - K-12 Level Assessment 16 \
\small (Answers are ordered exactly as in the question paper.)
\end{center}

\section*{Multiple Choice}
\begin{enumerate}[label=\arabic*.]
\item \textbf{Answer:} D
\\ \textit{Options:} A) what one sees when looking at bank literature; B) the nominal interest rate divided by the inflation rate; C) the nominal interest rate plus the anticipated inflation rate; D) the nominal interest rate minus anticipated inflation
\\ \textit{Note:} The real interest rate measures the true cost of borrowing and the true yield on savings by adjusting the nominal interest rate for the effects of anticipated inflation, reflecting the actual purchasing power of money over time.
\item \textbf{Answer:} C
\\ \textit{Options:} A) are a thing of the past.; B) are very severe depressions.; C) are marked by a sustained decline in output.; D) are regular occurrences in capitalist economies.
\\ \textit{Note:} Recessions are characterized by a sustained decline in economic output, typically measured by a decrease in gross domestic product (GDP), indicating a contraction in economic activity.
\item \textbf{Answer:} D
\\ \textit{Options:} A) U.S. GDP falls and Brazil's GDP falls.; B) U.S. GDP rises and Brazil's GDP falls.; C) U.S. GDP falls and Brazil's GDP remains constant.; D) U.S. GDP falls and Brazil's GDP rises.
\\ \textit{Note:} The correct answer is based on the fact that GDP measures the total value of goods and services produced within a country's borders, so when an American firm relocates its manufacturing to Brazil, the production previously counted in U.S. GDP is lost, while the new production in Brazil contributes to an increase in Brazil's GDP.
\item \textbf{Answer:} D
\\ \textit{Options:} A) the FED should actively conduct monetary policy.; B) changes in the money supply do not have significant effects.; C) fiscal policy is the preferred way of shifting the aggregate demand curve.; D) the FED should allow the money supply to grow at a constant rate.
\\ \textit{Note:} Monetarist theory posits that controlling the growth rate of the money supply is essential for maintaining economic stability and preventing inflation, advocating for a consistent, predictable increase in the money supply.
\item \textbf{Answer:} D
\\ \textit{Options:} A) Imports fall.; B) Income is transferred from steel consumers to domestic steel producers.; C) Income is transferred from foreign steel producers to domestic steel producers.; D) Allocative efficiency is improved.
\\ \textit{Note:} The removal of a protective tariff on imported steel allows for increased competition and lower prices, leading to a more efficient allocation of resources as consumers benefit from greater choices and producers are incentivized to innovate and improve efficiency.
\end{enumerate}

\section*{True / False}
\begin{enumerate}[label=\arabic*.]
\setcounter{enumi}{5}
\item \textbf{Answer:} True
\\ \textit{Options:} A) True; B) False
\\ \textit{Note:} The velocity of money measures how frequently a dollar is spent within a given period, reflecting the rate at which money circulates in the economy.
\item \textbf{Answer:} True
\\ \textit{Options:} A) True; B) False
\\ \textit{Note:} Higher government funding of research on clean energy supplies can lead to faster rates of economic growth by fostering innovation, creating new industries and jobs, and reducing energy costs, which collectively stimulate economic activity.
\item \textbf{Answer:} False
\\ \textit{Options:} A) True; B) False
\\ \textit{Note:} M$_{2}$ is not the least liquid measure of money because it includes components like savings accounts and time deposits, which are more liquid than M$_{3}$, the next broader measure of the money supply.
\item \textbf{Answer:} False
\\ \textit{Options:} A) True; B) False
\\ \textit{Note:} The statement is false because if the euro's value increases against the U.S. dollar, it means that the dollar has actually depreciated, not appreciated.
\item \textbf{Answer:} True
\\ \textit{Options:} A) True; B) False
\\ \textit{Note:} The statement is true because the shutdown price represents the lowest price at which a firm can cover its variable costs, which occurs at the minimum point of the average variable cost curve.
\end{enumerate}

\section*{Long Form Response}
\begin{enumerate}[label=\arabic*.]
\setcounter{enumi}{10}
\item \textbf{Answer:} Government subsidies can effectively address positive externalities in the production of private goods by providing financial support to firms or their customers. When a firm produces goods that generate benefits for society, such as education or renewable energy, these benefits are not fully reflected in the market price. By subsidizing the firm, the government lowers production costs, encouraging the firm to produce more of these beneficial goods. Alternatively, subsidizing customers can make these goods more affordable, increasing their consumption. Both approaches help align private incentives with social benefits, leading to a more efficient market outcome where the positive externalities are more fully realized.
\\ \textit{Note:} correct because it illustrates how government subsidies can lower production costs or consumer prices for goods that generate positive externalities, thereby incentivizing increased production and consumption of these socially beneficial goods, which leads to a more efficient market outcome.
\item \textbf{Answer:} Government intervention in a polluting steel company can significantly affect market equilibrium by reducing the company's output and increasing its prices. When the government imposes regulations, such as stricter environmental standards or pollution taxes, the cost of production for the steel company rises, leading to a decrease in the amount of steel produced. As a result, the supply of steel in the market declines. With less steel available, the price of steel tends to increase, as consumers compete for the limited supply. This change can create implications for consumers, who may face higher prices, and for the overall economy, as production levels adjust to meet new regulatory requirements.
\\ \textit{Note:} correct because government intervention increases production costs for the steel company, leading to a decrease in supply, which subsequently raises prices and alters market equilibrium.
\end{enumerate}

\vfill
\noindent\textit{End of Answer Key}
\end{document}